% ROZDZIAŁ 14 — samodzielny rozdział do projektu teoria_miary
\chapter{Przestrzenie $L^p$ i ich podstawowe własności}
\label{ch:przestrzenie-lp}

\section{Dlaczego utożsamiamy funkcje równe prawie wszędzie?}
Niech $(X,\mathcal F,\mu)$ będzie dowolną przestrzenią z miarą.
Nie zakładamy, że $\mu(X)<\infty$, ani że $\mu$ jest $\sigma$-skończona.
Dla $1\le p<\infty$ rozważmy mierzalne funkcje
$f:X\to\mathbb C$ spełniające
$\int_X|f(x)|^p\,\mathrm d\mu(x)<\infty$.
Naturalna wielkość
$(\int_X|f|^p\,\mathrm d\mu(x))^{1/p}$
może znikać dla funkcji nieidentycznie zerowej:
wystarczy, że $f$ różni się od zera jedynie na zbiorze
miary zero. Aby otrzymać normę, trzeba przejść
od funkcji do ich klas równoważności.

\begin{definicja}[Równość prawie wszędzie i przestrzeń $L^p$]
\label{def:lp-14}
Niech $(X,\mathcal F,\mu)$ będzie przestrzenią z miarą.
Dwie mierzalne funkcje zespolone $f,g$ na $X$
uznajemy za równoważne, jeśli $\mu(\{x:f(x)\ne g(x)\})=0$.
Dla $1\le p<\infty$ przez $L^p(X,\mathcal F,\mu)$, lub
krótko $L^p(\mu)$, rozumiemy przestrzeń klas
równoważności mierzalnych funkcji $f:X\to\mathbb C$,
dla których $\int_X|f(x)|^p\,\mathrm d\mu(x)<\infty$.
Wartość wyrażenia $\int_X|f|^p\,\mathrm d\mu(x)$
nie zależy od przedstawiciela klasy. Kładziemy
\begin{equation}\label{eq:p-norm-14}
\|f\|_p:=\left(\int_X|f(x)|^p\,\mathrm d\mu(x)\right)^{1/p}.
\end{equation}
Dla funkcji rzeczywistych używamy oznaczenia $L^p(\mu;\mathbb R)$.
\end{definicja}

\begin{definicja}[Kres górny istotny i przestrzeń $L^\infty$]
\label{def:linf-14}
Niech $(X,\mathcal F,\mu)$ będzie przestrzenią z miarą.
Dla mierzalnej funkcji $f:X\to\mathbb C$ określamy
\[
\|f\|_\infty:=\inf\{M\ge0:\mu(\{|f|>M\})=0\},
\]
z umową $\inf\varnothing=\infty$.
Przestrzeń $L^\infty(\mu)$ jest przestrzenią klas
funkcji o skończonej wartości $\|f\|_\infty$.
Liczbę tę nazywamy \emph{kresem górnym istotnym} funkcji $|f|$.
\end{definicja}

\begin{lemat}[Znaczenie kresu istotnego]
\label{lem:esssup-14}
Niech $(X,\mathcal F,\mu)$ będzie przestrzenią z miarą.
Jeśli $f:X\to\mathbb C$ jest mierzalna
i $\|f\|_\infty<\infty$,
to $|f(x)|\le\|f\|_\infty$ dla prawie wszystkich $x$.
Ponadto $\|f\|_\infty=0$ wtedy i tylko wtedy, gdy
$f=0$ prawie wszędzie.
\end{lemat}
\begin{proof}
Niech $M=\|f\|_\infty$. Dla każdego $n$ istnieje
$M_n<M+1/n$ spełniające $\mu(\{|f|>M_n\})=0$.
Poza przeliczalną sumą tych zbiorów zachodzi
$|f|\le M_n<M+1/n$ dla wszystkich $n$, więc $|f|\le M$.
W przypadku $M=0$ otrzymujemy $f=0$ poza zbiorem zerowym.
Implikacja odwrotna wynika wprost z definicji.
\end{proof}

\begin{uwaga}[Przedstawiciele a klasy]
Zapis $f(x)$ dla $f\in L^p(\mu)$ jest skrótem:
ściśle rzecz biorąc $f$ jest klasą równoważności, więc
wartość w pojedynczym punkcie nie musi być określona.
Ilekroć piszemy $f(x)$, wybieramy przedstawiciela;
tezy o wartościach punktowych formułujemy prawie wszędzie.
Operacje algebraiczne na klasach są poprawnie określone,
ponieważ suma skończonej liczby zbiorów zerowych jest zerowa.
\end{uwaga}

\section{Elementarna nierówność Younga}

\begin{lemat}[Young dla iloczynu dwóch liczb]
\label{lem:young-numbers-14}
Niech $1<p,q<\infty$ oraz $1/p+1/q=1$.
Dla $a,b\ge0$ zachodzi
\begin{equation}\label{eq:young-14}
ab\le\frac{a^p}{p}+\frac{b^q}{q}.
\end{equation}
Równość zachodzi dokładnie wtedy, gdy $a^p=b^q$.
\end{lemat}
\begin{proof}
Dla $b>0$ funkcja $\varphi(a)=a^p/p-ab+b^q/q$
ma pochodną $\varphi'(a)=a^{p-1}-b$.
Jej minimum na $[0,\infty)$ wypada w punkcie
$a=b^{1/(p-1)}=b^{q-1}$, gdzie
$\varphi(a)=0$, ponieważ $(q-1)p=q$.
Dla $b=0$ nierówność jest natychmiastowa;
równość zachodzi wtedy jedynie przy $a=0$.
\end{proof}

\begin{uwaga}
Wykładniki $p$ i $q$ w lemacie nazywamy
\emph{sprzężonymi wykładnikami Höldera}.
Dla $p=1$ przyjmujemy formalnie $q=\infty$ i odwrotnie.
Przypadki brzegowe rozpatrujemy osobno.
\end{uwaga}

\section{Nierówność Höldera}
\begin{twierdzenie}[Hölder]
\label{thm:holder-14}
Niech $(X,\mathcal F,\mu)$ będzie przestrzenią z miarą,
$1\le p,q\le\infty$ i $1/p+1/q=1$, przy czym $1/\infty=0$.
Jeśli $f\in L^p(\mu)$ i $g\in L^q(\mu)$,
to $fg\in L^1(\mu)$ oraz
\begin{equation}\label{eq:holder-14}
\int_X|f(x)g(x)|\,\mathrm d\mu(x)
\le\|f\|_p\|g\|_q.
\end{equation}
\end{twierdzenie}
\begin{proof}
Załóżmy najpierw $1<p,q<\infty$.
Jeśli $\|f\|_p=0$ lub $\|g\|_q=0$,
to odpowiedni czynnik znika prawie wszędzie,
a więc teza jest oczywista. W przeciwnym przypadku
kładziemy $F=|f|/\|f\|_p$ oraz $G=|g|/\|g\|_q$.
Zachodzi $\int_XF^p\,\mathrm d\mu(x)=
\int_XG^q\,\mathrm d\mu(x)=1$.
Z lematu~\ref{lem:young-numbers-14} mamy punktowo
$FG\le F^p/p+G^q/q$; całkowanie daje
$\int_XFG\,\mathrm d\mu(x)\le1/p+1/q=1$.
Po pomnożeniu przez $\|f\|_p\|g\|_q$ otrzymujemy
\eqref{eq:holder-14} i zarazem całkowalność iloczynu.

Gdy $p=1$ i $q=\infty$, poza zbiorem zerowym zachodzi
$|fg|\le\|g\|_\infty|f|$ na mocy
lematu~\ref{lem:esssup-14}. Całkowanie daje nierówność;
przypadek $p=\infty$, $q=1$ jest symetryczny.
\end{proof}

\begin{wniosek}[Uogólniony Hölder]
\label{cor:generalized-holder-14}
Niech $(X,\mathcal F,\mu)$ będzie przestrzenią z miarą,
$m\ge2$, $1\le p_j\le\infty$ oraz
$\sum_{j=1}^m1/p_j=1$. Jeżeli $f_j\in L^{p_j}(\mu)$,
to $\prod_{j=1}^m f_j\in L^1(\mu)$ i
$\int_X\prod_{j=1}^m|f_j(x)|\,\mathrm d\mu(x)
\le\prod_{j=1}^m\|f_j\|_{p_j}$.
\end{wniosek}
\begin{proof}
Czynniki z $p_j=\infty$ są ograniczone prawie wszędzie;
można je oszacować przez $\|f_j\|_\infty$ i usunąć
z rozważań. Jeśli pozostał tylko jeden czynnik,
jego wykładnik musi być równy $1$ i wynik jest oczywisty.
Dla większej liczby czynników dowodzimy twierdzenia indukcyjnie.
Niech $p_m<\infty$ i $1/r=\sum_{j<m}1/p_j=1-1/p_m$.
Jeżeli $r=\infty$, wszystkie pierwsze czynniki
są ograniczone i dostajemy tezę bezpośrednio.
Jeśli $1<r<\infty$, to $p_j\ge r$ dla $j<m$,
a liczby $s_j=p_j/r\ge1$ spełniają
$\sum_{j<m}1/s_j=1$. Stosujemy krok indukcyjny
do funkcji $|f_j|^r\in L^{s_j}$:
\[
\int_X\prod_{j<m}|f_j(x)|^r\,\mathrm d\mu(x)
\le\prod_{j<m}\left(\int_X|f_j(x)|^{p_j}\,
\mathrm d\mu(x)\right)^{r/p_j}.
\]
Po wyciągnięciu pierwiastka stopnia $r$ uzyskujemy
$\|\prod_{j<m}f_j\|_r\le\prod_{j<m}\|f_j\|_{p_j}$.
Ostatnie zastosowanie Höldera do tego iloczynu i $f_m$
ze sprzężonymi wykładnikami $r,p_m$ kończy dowód.
\end{proof}

\begin{przyklad}[Oszacowanie iloczynu]
Na $X=(0,1)$ z miarą Lebesgue'a rozważmy
$f(x)=x^{-1/4}$ i $g(x)=x^{-1/3}$.
Wtedy $f\in L^2(0,1)$, $g\in L^2(0,1)$,
ponieważ $\int_{(0,1)}x^{-1/2}\,\mathrm d\lambda_1(x)<\infty$
i $\int_{(0,1)}x^{-2/3}\,\mathrm d\lambda_1(x)<\infty$.
Z Höldera wynika, że $fg=x^{-7/12}\in L^1(0,1)$.
Można to również sprawdzić bezpośrednio.
\end{przyklad}

\section{Nierówność Minkowskiego i struktura normowana}
\begin{twierdzenie}[Minkowski]
\label{thm:minkowski-14}
Dla każdej przestrzeni z miarą $(X,\mathcal F,\mu)$,
każdego $1\le p\le\infty$ oraz $f,g\in L^p(\mu)$
zachodzi $f+g\in L^p(\mu)$ oraz
\begin{equation}\label{eq:minkowski-14}
\|f+g\|_p\le\|f\|_p+\|g\|_p.
\end{equation}
\end{twierdzenie}
\begin{proof}
Dla $p=1$ z nierówności trójkąta dla wartości
bezwzględnej i liniowości całki otrzymujemy
$\int_X|f+g|\,\mathrm d\mu(x)\le\int_X|f|\,\mathrm d\mu(x)+
\int_X|g|\,\mathrm d\mu(x)$.
Dla $p=\infty$ z
$|f+g|\le\|f\|_\infty+\|g\|_\infty$ prawie wszędzie
wynika teza.

Niech $1<p<\infty$, a $q=p/(p-1)$.
Na początek mamy
$|f+g|^p\le 2^{p-1}(|f|^p+|g|^p)$,
co wynika z wypukłości $t\mapsto t^p$ na $[0,\infty)$.
Ta nierówność zapewnia $f+g\in L^p$, zanim zastosujemy Höldera.
Następnie dla $s=f+g$ korzystamy z
$|s|^p\le |f||s|^{p-1}+|g||s|^{p-1}$.
Ponieważ $(p-1)q=p$, funkcja $|s|^{p-1}$ należy
do $L^q$ i $\||s|^{p-1}\|_q=\|s\|_p^{p-1}$.
Hölder daje
\[
\|s\|_p^p
\le(\|f\|_p+\|g\|_p)\|s\|_p^{p-1}.
\]
Gdy $\|s\|_p=0$, teza jest natychmiastowa;
w pozostałym przypadku dzielimy przez $\|s\|_p^{p-1}$.
\end{proof}

\begin{wniosek}[Przestrzeń normowana]
\label{cor:lp-normed-14}
Niech $(X,\mathcal F,\mu)$ będzie przestrzenią z miarą.
Dla $1\le p\le\infty$ wzory z
definicji~\ref{def:lp-14} i \ref{def:linf-14}
określają normę na zespolonej przestrzeni liniowej $L^p(\mu)$.
\end{wniosek}
\begin{proof}
Nieujemność jest oczywista; rozdzielność wynika
dla $p<\infty$ z faktu, że całka funkcji nieujemnej
równa zeru oznacza znikanie jej prawie wszędzie,
a dla $p=\infty$ z lematu~\ref{lem:esssup-14}.
Jednorodność $\|cf\|_p=|c|\|f\|_p$ wynika
z jednorodności całki lub kresu istotnego.
Nierówność trójkąta jest treścią twierdzenia
\ref{thm:minkowski-14}. Operacje na klasach
są niezależne od wyboru przedstawiciela.
\end{proof}

\section{Twierdzenie Riesza--Fischera: zupełność przestrzeni $L^p$}

\begin{twierdzenie}[Kryterium zbieżności szeregów w $L^p$]
\label{thm:series-lp-14}
Niech $(X,\mathcal F,\mu)$ będzie przestrzenią z miarą,
$1\le p\le\infty$ i $(u_n)_{n\ge1}\subset L^p(\mu)$.
Jeżeli $\sum_{n=1}^\infty\|u_n\|_p<\infty$,
to istnieje zbiór mierzalny $N$ o $\mu(N)=0$,
poza którym szereg $\sum_{n=1}^\infty u_n(x)$
jest bezwzględnie zbieżny. Jego suma określa element $u\in L^p(\mu)$
i $\sum_{n=1}^N u_n\to u$ w normie $L^p$.
\end{twierdzenie}
\begin{proof}
Dla $p<\infty$ kładziemy $v_N(x)=\sum_{n=1}^N|u_n(x)|$.
Nierówność Minkowskiego daje
$\|v_N\|_p\le\sum_{n=1}^N\|u_n\|_p\le M$,
gdzie $M=\sum_{n\ge1}\|u_n\|_p<\infty$.
Ponieważ $v_N\uparrow v:=\sum_{n\ge1}|u_n|$,
twierdzenie Beppo Leviego zastosowane do $v_N^p$ daje
$\int_Xv^p\,\mathrm d\mu(x)\le M^p$.
W konsekwencji $v<\infty$ prawie wszędzie.
Dla $x$ poza tym zbiorem wyjątkowym szereg
$\sum_nu_n(x)$ zbiega bezwzględnie do $u(x)$.
Wszystkie sumy częściowe są mierzalne, więc $u$
można określić jako funkcję mierzalną poza zbiorem
wyjątkowym i równą zero na nim. Mamy $|u|\le v$,
zatem $u\in L^p$.

Dla ogona $u-\sum_{n\le N}u_n$ mamy punktowo
$|u-\sum_{n\le N}u_n|\le\sum_{n>N}|u_n|$.
Stosując Minkowskiego do skończonych ogonów
i Fatou po przejściu z końcem ogona do nieskończoności,
otrzymujemy
$\|u-\sum_{n\le N}u_n\|_p\le\sum_{n>N}\|u_n\|_p\to0$.

Dla $p=\infty$ wybieramy przedstawicieli $u_n$
i jeden wspólny zbiór zerowy $N$, poza którym
$|u_n(x)|\le\|u_n\|_\infty$ dla każdego $n$.
Na $X\setminus N$ szereg zbiega bezwzględnie
oraz jednostajnie, gdyż suma $\sum_n\|u_n\|_\infty$
jest skończona. Otrzymujemy mierzalną sumę $u$,
$\|u\|_\infty\le M$ i
$\|u-\sum_{n\le N}u_n\|_\infty\le\sum_{n>N}\|u_n\|_\infty$.
\end{proof}

\begin{twierdzenie}[Riesz--Fischer: zupełność $L^p$]
\label{thm:riesz-fischer-14}
Dla każdej przestrzeni z miarą $(X,\mathcal F,\mu)$
i każdego $1\le p\le\infty$ przestrzeń normowana
$L^p(\mu)$ jest zupełna, a więc jest przestrzenią Banacha.
Dla $p=2$ jest ponadto przestrzenią Hilberta
z iloczynem skalarnym
$\langle f,g\rangle:=\int_X f(x)\overline{g(x)}\,\mathrm d\mu(x)$.
\end{twierdzenie}
\begin{proof}
Niech $(f_n)$ będzie ciągiem Cauchy'ego w $L^p$.
Indukcyjnie wybieramy rosnący ciąg indeksów $(n_k)$
taki, że $\|f_{n_{k+1}}-f_{n_k}\|_p\le2^{-k}$.
Wówczas szereg $\sum_{k\ge1}(f_{n_{k+1}}-f_{n_k})$
ma sumę norm mniejszą niż $\sum_k2^{-k}<\infty$.
Na mocy twierdzenia~\ref{thm:series-lp-14}
zbiega więc w $L^p$ do pewnego elementu $w$,
a zatem $f_{n_k}\to f:=f_{n_1}+w$ w $L^p$.
Aby wykazać zbieżność całego ciągu, ustalamy
$\varepsilon>0$ i wybieramy $N$ tak, aby
$\|f_m-f_n\|_p<\varepsilon/2$ dla $m,n\ge N$.
Wybieramy także $k$ z $n_k\ge N$ i
$\|f_{n_k}-f\|_p<\varepsilon/2$.
Wtedy dla $n\ge N$ nierówność trójkąta daje
$\|f_n-f\|_p\le\|f_n-f_{n_k}\|_p+
\|f_{n_k}-f\|_p<\varepsilon$.

Dla $p=2$ Hölder gwarantuje, że podana całka
jest skończona. Jest ona liniowa w pierwszym
argumencie, antyliniowa w drugim, hermitowska
i dodatnio określona na klasach. Wytworzona przez nią
norma to $\|f\|_2$, a udowodniona już zupełność
oznacza zupełność Hilberta.
\end{proof}

\begin{uwaga}[Zbieżność punktowa w dowodzie]
Istotnym elementem dowodu jest istnienie
\emph{podciągu} zbieżnego prawie wszędzie.
Nie wynika stąd, że dowolny ciąg zbieżny w $L^p$
zbiega punktowo prawie wszędzie; zobaczymy niżej
kontrprzykład związany z rodzajami zbieżności.
\end{uwaga}

\section{Nierówność Jensena i jej znaczenie}
\begin{twierdzenie}[Jensen dla miary probabilistycznej]
\label{thm:jensen-14}
Niech $(X,\mathcal F,P)$ będzie przestrzenią
probabilistyczną, $f\in L^1(P;\mathbb R)$,
a $\varphi:\mathbb R\to\mathbb R$ funkcją wypukłą.
Wtedy $\varphi\circ f$ jest mierzalna, jej część ujemna
jest całkowalna, a więc całka
$\int_X\varphi(f(x))\,\mathrm dP(x)\in(-\infty,\infty]$
jest dobrze określona i
\begin{equation}\label{eq:jensen-14}
\varphi\left(\int_Xf(x)\,\mathrm dP(x)\right)
\le\int_X\varphi(f(x))\,\mathrm dP(x).
\end{equation}
\end{twierdzenie}
\begin{proof}
Niech $m=\int_Xf\,\mathrm dP(x)$.
Funkcja wypukła skończona na całej prostej jest ciągła,
a więc $\varphi\circ f$ jest mierzalna.
Dla funkcji wypukłej na $\mathbb R$ w każdym punkcie $m$
istnieje prosta podpierająca: jest liczba $a\in\mathbb R$
taka, że $\varphi(t)\ge\varphi(m)+a(t-m)$
dla wszystkich $t\in\mathbb R$.
Istnienie takiego $a$ wynika z nierówności między
lewym i prawym nachyleniem siecznych funkcji wypukłej;
można wybrać dowolną liczbę między jej lewą
i prawą pochodną jednostronną w $m$.
Po podstawieniu $t=f(x)$ dostajemy
$\varphi(f(x))\ge\varphi(m)+a(f(x)-m)$.
Prawa strona jest całkowalna, więc część ujemna
$\varphi(f)$ ma skończoną całkę. Możemy całkować
nierówność, otrzymując
$\int_X\varphi(f)\,\mathrm dP(x)
\ge\varphi(m)+a(\int_Xf\,\mathrm dP(x)-m)=\varphi(m)$.
\end{proof}

\begin{przyklad}[Wartość średnia a momenty]
Dla $1\le p<\infty$ funkcja $\varphi(t)=|t|^p$
jest wypukła. Gdy $P(X)=1$ i $f\in L^p(P)$,
Jensen daje
$|\int_X f\,\mathrm dP(x)|^p\le
\int_X|f(x)|^p\,\mathrm dP(x)$ dla $f$ rzeczywistej.
Dla $f$ zespolonej analogiczna nierówność wynika
z nierówności trójkąta dla całki i Höldera
z funkcją stałą $1$ (należącą do $L^q(P)$).
\end{przyklad}

\section{Zależności między przestrzeniami $L^p$}

\begin{twierdzenie}[Włożenia przy skończonej mierze]
\label{thm:lp-embedding-finite-14}
Niech $(X,\mathcal F,\mu)$ będzie przestrzenią z miarą,
$\mu(X)<\infty$ i $1\le p<q\le\infty$.
Wtedy $L^q(\mu)\subseteq L^p(\mu)$ oraz
\begin{equation}\label{eq:finite-embedding-14}
\|f\|_p\le\mu(X)^{\,1/p-1/q}\|f\|_q
\quad(f\in L^q(\mu)),
\end{equation}
przy umowie $1/\infty=0$ i naturalnej interpretacji
w przypadku $\mu(X)=0$.
\end{twierdzenie}
\begin{proof}
Jeśli $\mu(X)=0$, wszystkie przestrzenie są
zerowe i teza jest natychmiastowa. Załóżmy więc
$0<\mu(X)<\infty$.
Dla $q<\infty$ stosujemy Höldera do
$|f|^p\cdot1$ z wykładnikami $q/p$ i $q/(q-p)$:
\[
\int_X|f|^p\,\mathrm d\mu(x)
\le\left(\int_X|f|^q\,\mathrm d\mu(x)\right)^{p/q}
\mu(X)^{1-p/q}.
\]
Podnosimy obie strony do potęgi $1/p$.
Dla $q=\infty$ nierówność
$|f|^p\le\|f\|_\infty^p$ prawie wszędzie
i całkowanie dają wynik.
\end{proof}

\begin{przyklad}[Na nieskończonej mierze włożenia zawodzą]
Na $(1,\infty)$ z miarą Lebesgue'a
funkcja $f(x)=x^{-\alpha}$ spełnia
$f\in L^p$ dokładnie wtedy, gdy $\alpha p>1$.
Na $(0,1)$ funkcja $g(x)=x^{-\beta}$
należy do $L^p$ dokładnie wtedy, gdy $\beta p<1$.
Dla $1\le p<q<\infty$ wybór
$1/q<\alpha\le1/p$ daje na $(1,\infty)$
funkcję z $L^q\setminus L^p$,
a wybór $1/q\le\beta<1/p$ daje na $(0,1)$
funkcję z $L^p\setminus L^q$.
Po przedłużeniu przez zero odpowiednie przykłady
działają na całej $\mathbb R$.
Zatem na przestrzeni nieskończonej miary
w ogólności nie zachodzi żadna z inkluzji między
różnymi $L^p$.
\end{przyklad}

\begin{twierdzenie}[Zbieżność w $L^p$ a zbieżność według miary]
\label{thm:lp-to-measure-14}
Niech $(X,\mathcal F,\mu)$ będzie przestrzenią z miarą,
$1\le p<\infty$ i niech $f_n,f\in L^p(\mu)$.
Jeżeli $\|f_n-f\|_p\to0$, to $f_n\to f$
według miary $\mu$, tj. dla każdego $\varepsilon>0$
$\mu(\{|f_n-f|>\varepsilon\})\to0$.
Dla $p=\infty$ zachodzi silniejsze stwierdzenie:
$\|f_n-f\|_\infty\to0$ implikuje
$\mu(\{|f_n-f|>\varepsilon\})=0$ dla dostatecznie
dużych $n$.
\end{twierdzenie}
\begin{proof}
Dla $p<\infty$ nierówność Czebyszewa daje
$\varepsilon^p\mu(\{|f_n-f|>\varepsilon\})
\le\int_X|f_n-f|^p\,\mathrm d\mu(x)=\|f_n-f\|_p^p$.
Prawa strona zbiega do zera.
Dla $p=\infty$ wybieramy $n$ tak duże, że
$\|f_n-f\|_\infty<\varepsilon$ i korzystamy
z definicji kresu istotnego.
\end{proof}

\begin{przyklad}[Zbieżność według miary nie wystarcza]
Na $(0,1)$ z miarą Lebesgue'a połóżmy
$f_n(x)=n^{1/p}\chi_{(0,1/n)}(x)$ dla $p<\infty$.
Wtedy $f_n\to0$ według miary i prawie wszędzie,
ale $\|f_n\|_p=1$ dla każdego $n$.
Zbieżność według miary nie pociąga więc zbieżności
w normie, nawet na przestrzeni probabilistycznej.
\end{przyklad}

\section{Interpolacja pomiędzy normami $L^p$}

\begin{twierdzenie}[Nierówność interpolacyjna]
\label{thm:lp-interpolation-14}
Niech $(X,\mathcal F,\mu)$ będzie przestrzenią z miarą,
$1\le p<r<q\le\infty$, i niech $\theta\in(0,1)$
będzie jedyną liczbą spełniającą
$1/r=\theta/p+(1-\theta)/q$, przy czym $1/\infty=0$.
Dla każdej funkcji $f\in L^p(\mu)\cap L^q(\mu)$
zachodzi $f\in L^r(\mu)$ oraz
\begin{equation}\label{eq:lp-interpolation-14}
\|f\|_r\le\|f\|_p^{\theta}\|f\|_q^{1-\theta}.
\end{equation}
W szczególności nie zakładamy $\mu(X)<\infty$.
\end{twierdzenie}
\begin{proof}
Najpierw załóżmy $q<\infty$ i obie normy po prawej
dodatnie. Wybieramy sprzężone wykładniki
$a=p/(\theta r)>1$ oraz $b=q/((1-\theta)r)>1$;
warunek na $\theta$ jest równoważny $1/a+1/b=1$.
Z twierdzenia Höldera dla czynników
$|f|^{\theta r}$ i $|f|^{(1-\theta)r}$ otrzymujemy
\[
\int_X|f(x)|^r\,\mathrm d\mu(x)
\le\left(\int_X|f(x)|^p\,\mathrm d\mu(x)\right)^{\theta r/p}
\left(\int_X|f(x)|^q\,\mathrm d\mu(x)\right)^{(1-\theta)r/q}.
\]
Podnosząc obie strony do potęgi $1/r$, dostajemy tezę.
Gdy $q=\infty$, warunek na $\theta$ daje
$\theta r=p$. Wówczas
$|f|^r\le\|f\|_\infty^{r-p}|f|^p$
prawie wszędzie, z czego po całkowaniu wynika
\eqref{eq:lp-interpolation-14}. Jeżeli któraś
z norm po prawej jest zerowa, funkcja znika
prawie wszędzie i nierówność jest oczywista.
\end{proof}

\begin{przyklad}[Interpolacja bez skończoności miary]
Na $X=\mathbb R^d$ założenie $f\in L^1\cap L^\infty$
gwarantuje $f\in L^r$ dla wszystkich $1<r<\infty$,
a dokładniej $\|f\|_r\le\|f\|_1^{1/r}
\|f\|_\infty^{1-1/r}$. Nie wolno jednak wnioskować
z samego $f\in L^q$ o przynależności do $L^p$
przy $p<q$ bez dodatkowych założeń o mierze.
\end{przyklad}

\section{Aproksymacja i gęstość wybranych klas funkcji}

\begin{twierdzenie}[Gęstość funkcji prostych]
\label{thm:simple-dense-lp-14}
Niech $(X,\mathcal F,\mu)$ będzie dowolną przestrzenią
z miarą i $1\le p<\infty$. Wtedy w $L^p(\mu)$ gęste są
klasy funkcji prostych, przyjmujących skończenie wiele
wartości i znikających poza mierzalnym zbiorem skończonej miary.
Dla $p=\infty$ funkcje proste są gęste w $L^\infty(\mu)$,
ale bez dodatkowych założeń nie można żądać, aby
znikały poza zbiorami skończonej miary.
\end{twierdzenie}
\begin{proof}
Najpierw niech $f\ge0$ i $f\in L^p$ dla $p<\infty$.
Dla $n\ge1$ kładziemy $E_n=\{f>1/n\}$.
Z nierówności Czebyszewa mamy
$\mu(E_n)\le n^p\int_Xf^p\,\mathrm d\mu(x)<\infty$.
Funkcje $f\chi_{E_n}\chi_{\{f\le n\}}$
zbiegają punktowo do $f$ poza zbiorem, na którym
$f=\infty$ (zerowym, gdy $f\in L^p$), a ich
różnice z $f$ są zdominowane przez $f$.
Twierdzenie Lebesgue'a o zbieżności zmajoryzowanej
zastosowane do $p$-tych potęg modułów różnic pokazuje,
że obcięcia zbiegają do $f$ w $L^p$.

Dla ustalonych $n$ obcięta funkcja jest ograniczona
przez $n$ i znika poza $E_n$ o skończonej mierze.
Aproksymujemy ją z dołu funkcjami prostymi
$o$ skoku wartości co najwyżej $1/m$; błąd
jest co najwyżej $1/m$ na $E_n$, a zatem jego
norma $L^p$ jest co najwyżej
$\mu(E_n)^{1/p}/m\to0$. Dla funkcji rzeczywistej
o dowolnym znaku stosujemy to rozumowanie
osobno do jej części dodatniej i ujemnej;
obie należą do $L^p(\mu)$, a różnica ich
aproksymacji ma skończony zbiór wartości
i znika poza sumą dwóch zbiorów
skończonej miary. Dla funkcji zespolonej
wykonujemy tę konstrukcję osobno dla części
rzeczywistej i urojonej, a następnie
korzystamy z nierówności Minkowskiego.

Dla $p=\infty$ każdą funkcję istotnie ograniczoną
można poza zbiorem zerowym aproksymować
jednostajnie funkcjami prostymi, przez podział
ograniczonego prostokąta wartości w $\mathbb C$
na kwadraty o dowolnie małych średnicach.
Jeżeli $\mu(X)=\infty$, funkcja $1\in L^\infty(\mu)$,
ale odległość w normie $L^\infty$ między nią
a każdą funkcją zerową poza zbiorem skończonej
miary wynosi co najmniej $1$.
\end{proof}

\begin{twierdzenie}[Gęstość funkcji gładkich w przestrzeni euklidesowej]
\label{thm:cc-dense-lp-14}
Dla $d\ge1$ oraz $1\le p<\infty$ zbiór
$C_c^\infty(\mathbb R^d)$ jest gęsty w
$L^p(\mathbb R^d,\mathcal L_d,\lambda_d)$.
Nie jest gęsty w $L^\infty(\mathbb R^d)$.
\end{twierdzenie}
\begin{proof}
Przypadek $1\le p<\infty$ został wykazany
w rozdziale~\ref{ch:regularnosc-funkcji},
w twierdzeniu~\ref{thm:cc-in-lp-10}:
korzysta on z regularności miary Lebesgue'a,
aproksymacji funkcji charakterystycznych zbiorów mierzalnych
funkcjami ciągłymi o zwartym nośniku,
a następnie z wygładzania przez mollifikatory.
Przytoczenie tego wyniku nie wykorzystuje
żadnego twierdzenia z późniejszych rozdziałów.

Dla $p=\infty$ rozważmy funkcję stałą $1$.
Dla dowolnej $\varphi\in C_c^\infty(\mathbb R^d)$
poza zwartym nośnikiem $\varphi$ zachodzi $\varphi=0$,
a zbiór ten ma dodatnią miarę. Stąd
$\|1-\varphi\|_\infty\ge1$, więc gęstość nie zachodzi.
\end{proof}

\begin{twierdzenie}[Ciągłość translacji w $L^p(\mathbb R^d)$]
\label{thm:translation-lp-14}
Niech $d\ge1$, $1\le p<\infty$ oraz
$f\in L^p(\mathbb R^d,\mathcal L_d,\lambda_d)$.
Dla $h\in\mathbb R^d$ kładziemy
$(\tau_hf)(x)=f(x-h)$.
Wówczas $\|\tau_hf-f\|_p\to0$ przy $h\to0$.
\end{twierdzenie}
\begin{proof}
Niezmienniczość miary Lebesgue'a względem translacji
daje $\|\tau_hg\|_p=\|g\|_p$.
Jeżeli $g\in C_c^\infty(\mathbb R^d)$, to jest
jednostajnie ciągła, a dla $|h|\le1$
nośniki $g(\cdot-h)$ oraz $g$ leżą we wspólnym
zbiorze zwartym $K$ o skończonej mierze.
Mamy zatem
$\|\tau_hg-g\|_p\le
\lambda_d(K)^{1/p}\sup_x|g(x-h)-g(x)|\to0$.

Dla dowolnego $f\in L^p$ i $\varepsilon>0$
wybieramy z gęstości $g\in C_c^\infty$ tak, że
$\|f-g\|_p<\varepsilon$.
Nierówność trójkąta prowadzi do
$\|\tau_hf-f\|_p\le
\|\tau_h(f-g)\|_p+\|\tau_hg-g\|_p+\|g-f\|_p
<2\varepsilon+\|\tau_hg-g\|_p$.
Biorąc granicę górną przy $h\to0$ i następnie
$\varepsilon\downarrow0$, otrzymujemy tezę.
\end{proof}

\begin{przyklad}[Dlaczego $p=\infty$ jest inny]
Dla $f=\chi_{(0,\infty)}$ na $\mathbb R$
$\|\tau_hf-f\|_\infty=1$ dla każdego $h\ne0$.
Zatem translacje nie są silnie ciągłe w
$L^\infty(\mathbb R)$.
\end{przyklad}

\section{Separowalność: kiedy istnieje przeliczalna rodzina gęsta?}

\begin{twierdzenie}[Separowalność $L^p$]
\label{thm:separable-lp-14}
Niech $(X,\mathcal F,\mu)$ będzie przestrzenią
z miarą $\sigma$-skończoną, przy czym $\mathcal F$
jest generowana przez przeliczalną rodzinę zbiorów.
Wówczas $L^p(\mu)$ jest separowalna dla
każdego $1\le p<\infty$.
\end{twierdzenie}
\begin{proof}
Wybieramy rosnące wyczerpanie $X_n\uparrow X$
zbiorami o skończonej mierze. Do przeliczalnej
rodziny generującej $\mathcal F$ dodajemy zbiory $X_n$
i oznaczamy przez $\mathcal A$ algebrę przez nią generowaną.
Ponieważ wszystkie operacje w algebrze są skończone,
$\mathcal A$ jest przeliczalna i $\sigma(\mathcal A)=\mathcal F$.

Na każdym $X_n$ rozważamy klasę zbiorów $E\in\mathcal F|_{X_n}$,
dla których dla każdego $\varepsilon>0$ istnieje
$A\in\mathcal A|_{X_n}$ z
$\mu(E\mathbin\triangle A)<\varepsilon$.
Klasa ta jest $\sigma$-algebrą: dla dwóch zbiorów
aproksymowanych zbiorami $A,B$ z algebry zachodzi
$\mu((E\cup F)\mathbin\triangle(A\cup B))\le
\mu(E\mathbin\triangle A)+\mu(F\mathbin\triangle B)$;
analogicznie traktujemy dopełnienia względem $X_n$.
Dla $E_k\uparrow E$ mamy
$\mu(E\setminus E_k)\to0$ dzięki
$\mu(X_n)<\infty$. Stąd przeliczalne sumy
aproksymujemy najpierw skończonymi, a następnie
zbiorami z algebry. Klasa zawiera $\mathcal A|_{X_n}$,
a więc jest całą $\mathcal F|_{X_n}$.

Z twierdzenia~\ref{thm:simple-dense-lp-14}
wystarczy aproksymować funkcje proste znikające
poza zbiorem skończonej miary. Ponieważ $X_n\uparrow X$, dla każdej
$f\in L^p(\mu)$ mamy
$\int_X |f|^p\chi_{X_n}\,\mathrm d\mu
\uparrow\int_X|f|^p\,\mathrm d\mu$
na mocy twierdzenia Beppo Leviego.
Stąd $\|f\chi_{X\setminus X_n}\|_p\to0$;
w szczególności zastąpienie funkcji prostej
przez jej iloczyn z $\chi_{X_n}$ ma dowolnie mały błąd. Następnie każdą
funkcję charakterystyczną na $X_n$ aproksymujemy
w $L^p$ przez funkcje charakterystyczne zbiorów z
$\mathcal A|_{X_n}$, bo
$\|\chi_E-\chi_A\|_p
=\mu(E\mathbin\triangle A)^{1/p}$.
Wreszcie współczynniki zespolone aproksymujemy
liczbami o wymiernych częściach rzeczywistych
i urojonych. Zbiór wszystkich takich skończonych
kombinacji funkcji charakterystycznych zbiorów z przeliczalnej
algebry (z obcięciem do wybranych $X_n$) jest
przeliczalny i gęsty w $L^p$.
\end{proof}

\begin{przyklad}[Dwie przeszkody dla separowalności]
Dla $X=[0,1]$ z miarą Lebesgue'a przestrzeń
$L^\infty(X)$ nie jest separowalna: funkcje
$\chi_{(0,t)}$, $t\in(0,1)$, są parami
oddalone w normie $L^\infty$ o $1$.
Zbiór nieprzeliczalny, którego różne elementy są
od siebie oddalone co najmniej o $1$, nie może
mieścić się w przestrzeni separowalnej.

Również dla $p<\infty$ istnieją przestrzenie
nieseparowalne: jeśli $I$ jest nieprzeliczalny,
a $\kappa$ jest miarą liczącą na $\mathcal P(I)$,
to funkcje $e_i=\chi_{\{i\}}\in L^p(\kappa)$
dla $i\in I$ mają parami odległość $2^{1/p}$.
Tutaj założenie o przeliczalnym generowaniu
$\sigma$-algebry zawodzi.
\end{przyklad}

\section{Co dalej: dualność i słaba zbieżność}
Dzięki zupełności możemy traktować $L^p$ jako
przestrzenie Banacha, a $L^2$ jako przestrzeń Hilberta.
Nie przesądza to jednak ich struktury dualnej.
W następnym rozdziale zbadamy, dla jakich
wykładników i przestrzeni z miarą każdy ciągły
funkcjonał liniowy na $L^p$ można zapisać jako
$f\mapsto\int_Xf(x)g(x)\,\mathrm d\mu(x)$.
Będziemy tam precyzyjnie odróżniać
przestrzeń dualną od samego zbioru funkcji reprezentujących,
a także omówimy refleksywność i zbieżność słabą.
